\documentclass[10pt]{article}
\usepackage[utf8]{inputenc}
\usepackage{TLCresume}
\usepackage[margin=1in, top=1in, bottom=1in]{geometry}

%====================
% CONTACT INFORMATION
%====================
\def\name{Nishit Panchal}
\def\phone{+91 9724286921}
\def\city{Ahmedabad, Gujarat, India}
\def\email{nishit1024@gmail.com}
\def\LinkedIn{nishit-panchal-346353145}
\def\role{Senior iOS Engineer}

\begin{document}
\pagenumbering{gobble}

%====================
% HEADER
%====================
\begin{center}
    {\Huge \textbf{\name}} \\
    \vspace{0.4em}
    {\color{highlight} \Large \role} \\
    \vspace{0.5em}
    \small{
        \href{mailto:\email}{\email} \ \textbullet \
        \href{https://www.linkedin.com/in/\LinkedIn}{linkedin.com/in/\LinkedIn} \ \textbullet \
        \city
    }
\end{center}

\vspace{1em}
\hrule
\vspace{1.5em}

%====================
% DATE & RECIPIENT
%====================
\noindent
Hiring Team \\
Apple Inc. \\

\vspace{1em}

%====================
% SALUTATION
%====================
\noindent Dear Hiring Team,

\vspace{1em}

%====================
% BODY
%====================
\noindent
I still remember the first time I held an iPhone and thought --- how does something this small feel this alive? That question has driven every line of code I have written for the past eight years.

\vspace{0.8em}

\noindent
I am Nishit, a Senior iOS Engineer with a deep specialisation in Apple's UI frameworks --- SwiftUI, UIKit, Core Animation, and Core Graphics. Over the course of my career I have taken 6 iOS applications from a blank Xcode project to the App Store, collectively serving more than 500,000 users. I have done this in fast-moving startup environments where there was no safety net --- just a clear vision, a small team, and the expectation that every release would be better than the last. That pressure built something in me: an obsession with craft, with performance, and with the feeling a user gets when an interface responds exactly the way their instincts expected it to.

\vspace{0.8em}

\noindent
The opportunity on your team speaks directly to that obsession. Building UI frameworks and experiences that surface machine learning in ways that feel effortless --- not technical --- is exactly the kind of challenge I have been moving toward. My work with ARKit and RealityKit taught me that the hardest part of ML-powered features is not the model. It is making the output feel like magic to someone who has no idea what is happening under the hood. That translation --- from computation to experience --- is where I do my best work.

\vspace{0.8em}

\noindent
I have also spent meaningful time with Xcode Instruments, using Time Profiler and Memory Graph Debugger to hunt down the kind of subtle performance issues that users feel before they can name them. A dropped frame, a brief stutter, a moment of hesitation --- these things matter enormously when you are building experiences that people return to every day. I believe that performance is a feature, and I have built my engineering habits around that belief.

\vspace{0.8em}

\noindent
Beyond the technical work, I have led cross-functional teams, mentored 14+ junior developers, and operated in fully remote environments where clear communication and ownership were non-negotiable. I know how to work with designers and product managers to turn ambiguous ideas into concrete, shippable increments --- and I know how to hold myself accountable to quality even when no one is checking.

\vspace{0.8em}

\noindent
Apple has always set the standard for what software can feel like. I would be proud to contribute to that standard --- and genuinely excited to work on experiences that sit at the frontier of what Apple Intelligence makes possible.

\vspace{0.8em}

\noindent
Thank you for taking the time to read this. I would love the opportunity to talk.

\vspace{1.5em}

%====================
% SIGN OFF
%====================
\noindent
Warm regards, \\
\vspace{0.3em}
{\bfseries Nishit Panchal} \\
\href{mailto:\email}{\email} \\
\href{https://www.linkedin.com/in/\LinkedIn}{linkedin.com/in/\LinkedIn}

\end{document}
